\documentclass[11pt,a4paper]{article}
\usepackage[T1]{fontenc}
\usepackage[utf8]{inputenc}
\usepackage{lmodern,amsmath,amssymb}
\usepackage[margin=25mm]{geometry}
\usepackage[hidelinks]{hyperref}
\hypersetup{pdftitle={The flowed bound in free field theory: the coupling cancels, and the bound is informative },pdfauthor={navstokgap research working notes}}
\newcommand{\tightlist}{\setlength{\itemsep}{0pt}\setlength{\parskip}{0pt}}
\setlength{\emergencystretch}{3em}
\setcounter{secnumdepth}{0}
\begin{document}
\hypertarget{the-flowed-bound-in-free-field-theory-the-coupling-cancels-and-the-bound-is-informative-only-at-the-confinement-scale}{%
\section{The flowed bound in free field theory: the coupling cancels,
and the bound is informative only at the confinement
scale}\label{the-flowed-bound-in-free-field-theory-the-coupling-cancels-and-the-bound-is-informative-only-at-the-confinement-scale}}

The two ground-state quantities of
\href{finiteness-half-flowed-susceptibility.md}{the finiteness note} are
computed in closed form for the free (abelian) theory with the flowed
magnetic energy density
\(\varphi_t=\tfrac12(b_t\!\cdot\!b_t-\langle\cdot\rangle)\):
\[\chi_{\varphi_t}=\frac{3g^4}{2048\,\pi^{3/2}\,t^{5/2}},\qquad
\big\langle|D_{\varphi_t}|^2\big\rangle=\frac{g^2}{128\,\pi^2\,t^3},\qquad
\frac{\langle|D|^2\rangle}{\chi}=\frac{16}{3\,g^2\sqrt{\pi t}},\] so
that the bound \(m\le\frac{g^2\hbar c}{2}\langle|D|^2\rangle/\chi\)
becomes
\[\boxed{\;m\ \le\ \frac{8}{3\sqrt\pi}\,\frac{\hbar c}{\sqrt t}\;=\;\frac{16\sqrt2}{3\sqrt\pi}\,\frac{\hbar c}{\sqrt{8t}}\;\approx\;4.26\,\frac{\hbar c}{\sqrt{8t}}\;}\]
with the coupling cancelling between numerator and denominator. This
corrects Section 5 of the finiteness note, which carried a factor
\(g^2\) from dimensional analysis alone: \(\chi\propto g^4\) and
\(\langle|D|^2\rangle\propto g^2\), so the ratio supplies \(g^{-2}\).
The result is confirmed independently by the spectral form of the same
bound: the smeared operator creates two photons of momenta \(\pm k\)
with weight \(\propto k^4e^{-4tk^2}dk\), and the weighted mean of their
energy \(2\hbar ck\) is exactly \(8\hbar c/(3\sqrt{\pi t})\). Two
consequences. The free bound tends to zero as the smearing grows, as it
must for \(m=0\), so the inequality retains content. And in the
interacting theory the weak-coupling regime is
\(\sqrt{8t}\ll\hbar c/(\text{confinement scale})\), where the bound
reads \(m\lesssim4.26\,\hbar c/\sqrt{8t}\) and is far weaker than the
expected gap: the bound becomes informative exactly at
\(\sqrt{8t}\sim\hbar c/m\), where its coefficient is no longer
perturbative. No perturbative computation can make the finiteness half
useful; what is needed is the value of one susceptibility at the
confinement scale. Constants explicit; nothing promoted.

\hypertarget{the-free-ground-state-with-explicit-constants}{%
\subsection{1. The free ground state with explicit
constants}\label{the-free-ground-state-with-explicit-constants}}

Take the abelian case of the Hamiltonian of
\href{finiteness-half-flowed-susceptibility.md}{the finiteness note} §1,
\[H=\int d^3x\;\Big[\frac{g^2c}{2\hbar}\,\Pi_i\Pi_i+\frac{\hbar c}{2g^2}\,b_ib_i\Big],
\qquad b_i=\varepsilon_{ijk}\partial_jA_k,
\qquad[A_i(x),\Pi_j(y)]=i\hbar\delta_{ij}\delta^3(x-y),\] in Coulomb
gauge \(\partial_iA_i=0\), two transverse polarizations. In Fourier
modes each transverse mode is an oscillator of mass \(M=\hbar/(g^2c)\)
and frequency \(\omega_k=ck\), since
\(\tfrac12M\omega_k^2=\hbar ck^2/(2g^2)\) matches the magnetic term. Its
ground-state width gives
\[\big\langle\hat A_i(k)\hat A_j(k')\big\rangle_0=(2\pi)^3\delta^3(k+k')\,P^{\rm T}_{ij}(k)\,\frac{\hbar}{2M\omega_k}
=(2\pi)^3\delta^3(k+k')\,P^{\rm T}_{ij}(k)\,\frac{g^2}{2k},\] with
\(P^{\rm T}_{ij}=\delta_{ij}-\hat k_i\hat k_j\). Dimensions:
\([\hat A]=L^2\), \([(2\pi)^3\delta^3]=L^3\), and \(g^2/2k\) is a
length.

\hypertarget{the-flow-is-the-heat-equation-here}{%
\subsection{2. The flow is the heat equation
here}\label{the-flow-is-the-heat-equation-here}}

For an abelian field the flow equation of
\href{finiteness-half-flowed-susceptibility.md}{the finiteness note} §4
reduces to \(\partial_sB_i=\Delta B_i-\partial_i(\partial_jB_j)\), and a
transverse initial field stays transverse, so
\[\hat B_i(t,k)=e^{-tk^2}\hat A_i(k),\qquad
\hat b_i(t,k)=i\varepsilon_{ijl}k_j\,e^{-tk^2}\hat A_l(k).\] The flow
time has dimension of length squared; the conventional smearing radius
is \(\sqrt{8t}\). Set
\(\varphi_t=\tfrac12\big(b_t\!\cdot\!b_t-\langle b_t\!\cdot\!b_t\rangle_0\big)\),
of dimension \(L^{-4}\).

\hypertarget{the-susceptibility}{%
\subsection{3. The susceptibility}\label{the-susceptibility}}

Write \(G_{ij}(z)=\langle b_{t,i}(z)b_{t,j}(0)\rangle_0\). Using
\(\varepsilon_{iab}\varepsilon_{jcd}k_ak_cP^{\rm T}_{bd}(k)=k^2P^{\rm T}_{ij}(k)\)
(the longitudinal part of \(P^{\rm T}\) drops because
\(\varepsilon_{iab}k_a\hat k_b=0\)),
\[\tilde G_{ij}(k)=k^2P^{\rm T}_{ij}(k)\,\frac{g^2}{2k}\,e^{-2tk^2}
=\frac{g^2k}{2}\,e^{-2tk^2}\,P^{\rm T}_{ij}(k).\] For a Gaussian field
Wick's theorem gives
\(\langle\varphi_t(z)\varphi_t(0)\rangle^{\rm c}_0=\tfrac12\sum_{ij}G_{ij}(z)^2\)
(two pairings), so by Parseval, with
\(\sum_{ij}[P^{\rm T}_{ij}]^2=\operatorname{tr}P^{\rm T}=2\),
\[\chi_{\varphi_t}=\int d^3z\,\langle\varphi_t\varphi_t\rangle^{\rm c}
=\frac12\int\frac{d^3k}{(2\pi)^3}\,\frac{g^4k^2}{2}\,e^{-4tk^2}
=\frac{g^4}{4}\cdot\frac1{2\pi^2}\int_0^\infty\!dk\,k^4e^{-4tk^2}.\]
With \(\int_0^\infty k^4e^{-\alpha k^2}dk=3\sqrt\pi/(8\alpha^{5/2})\)
and \(\alpha=4t\), so \((4t)^{5/2}=32t^{5/2}\),
\[\chi_{\varphi_t}=\frac{g^4}{4}\cdot\frac1{2\pi^2}\cdot\frac{3\sqrt\pi}{256\,t^{5/2}}
=\frac{3g^4}{2048\,\pi^{3/2}\,t^{5/2}} .\] Dimension:
\([\chi]=[\varphi]^2L^3=L^{-5}\) and \(t^{5/2}\sim L^5\), as required.

\hypertarget{the-gradient-term}{%
\subsection{4. The gradient term}\label{the-gradient-term}}

By Parseval
\(O=\int d^3y\,\varphi_t=\tfrac12\int\frac{d^3k}{(2\pi)^3}k^2e^{-2tk^2}|\hat A(k)|^2+\text{const}\),
so
\[D_i(x)=\frac{\delta O}{\delta A_i(x)}=\big(-\Delta\,e^{2t\Delta}A\big)_i(x)
=\big(\nabla\times b_{2t}\big)_i(x),\qquad
\hat D_i(k)=k^2e^{-2tk^2}\hat A_i(k),\] the curl of the magnetic field
at twice the flow time. Hence
\[\big\langle|D|^2\big\rangle_0=\int\frac{d^3k}{(2\pi)^3}k^4e^{-4tk^2}\cdot2\cdot\frac{g^2}{2k}
=g^2\cdot\frac1{2\pi^2}\int_0^\infty\!dk\,k^5e^{-4tk^2}
=\frac{g^2}{2\pi^2}\cdot\frac1{64\,t^3}=\frac{g^2}{128\,\pi^2\,t^3},\]
using \(\int_0^\infty k^5e^{-\alpha k^2}dk=1/\alpha^3\). Dimension:
\([\varphi]^2L^2=L^{-6}\) and \(t^3\sim L^6\), as required.

\hypertarget{the-bound-and-the-cancellation-of-the-coupling}{%
\subsection{5. The bound, and the cancellation of the
coupling}\label{the-bound-and-the-cancellation-of-the-coupling}}

\[\frac{\langle|D|^2\rangle}{\chi}
=\frac{g^2}{128\pi^2t^3}\cdot\frac{2048\,\pi^{3/2}t^{5/2}}{3g^4}
=\frac{16}{3\,g^2\sqrt{\pi t}},\]
\[m\ \le\ \frac{g^2\hbar c}{2}\cdot\frac{16}{3g^2\sqrt{\pi t}}
=\frac{8}{3\sqrt\pi}\,\frac{\hbar c}{\sqrt t}
=\frac{16\sqrt2}{3\sqrt\pi}\,\frac{\hbar c}{\sqrt{8t}}\approx4.26\,\frac{\hbar c}{\sqrt{8t}} .\]

\textbf{Correction.} Section 5 of
\href{finiteness-half-flowed-susceptibility.md}{the finiteness note}
inferred \(m\le Cg^2\hbar c/\sqrt{8t}\) from dimensional analysis,
assigning the pure numbers \(k_1,k_2\) no coupling dependence. The
computation shows \(\chi\propto g^4\) while
\(\langle|D|^2\rangle\propto g^2\), so the ratio carries \(g^{-2}\) and
the bound is of order \(g^0\). The corrected statement is that the
perturbative bound is \(4.26\,\hbar c/\sqrt{8t}\,[1+O(g^2)]\), with the
coupling entering only through the corrections. The \(t\)-dependence
\(1/\sqrt{8t}\) of the earlier note stands.

\hypertarget{independent-check-the-two-photon-sum-rule}{%
\subsection{6. Independent check: the two-photon sum
rule}\label{independent-check-the-two-photon-sum-rule}}

By the spectral reading of
\href{finiteness-half-flowed-susceptibility.md}{the finiteness note} §3,
the bound equals the mean excitation energy weighted by the spectral
weight of \(O\). In the free theory \(O\) is quadratic in \(\hat A\) and
creates two photons of momenta \(k\) and \(-k\), of total energy
\(2\hbar ck\), with weight proportional to
\(\big[k^2e^{-2tk^2}\big]^2\big[g^2/2k\big]^2\,k^2\,dk\propto k^4e^{-4tk^2}dk\)
after the phase-space factor. Therefore
\[\frac{\int_0^\infty 2\hbar ck\cdot k^4e^{-4tk^2}dk}{\int_0^\infty k^4e^{-4tk^2}dk}
=2\hbar c\,\frac{1/(64t^3)}{3\sqrt\pi/(256\,t^{5/2})}
=2\hbar c\cdot\frac{4}{3\sqrt{\pi t}}
=\frac{8}{3\sqrt\pi}\,\frac{\hbar c}{\sqrt t},\] the same number. The
agreement checks Proposition 1, the intensive form of Corollary 2 and
the Gaussian computation against each other, since the two routes share
no step beyond the definition of \(\varphi_t\).

\hypertarget{what-this-says-about-the-interacting-theory}{%
\subsection{7. What this says about the interacting
theory}\label{what-this-says-about-the-interacting-theory}}

\emph{The free limit is correct.} For \(m=0\) the bound must be able to
vanish, and it does, like \(1/\sqrt{8t}\): the inequality is not
vacuous, and the flow time is the only resolution scale in the free
theory.

\emph{Weak coupling proves nothing.} In the interacting theory the
perturbative regime at flow time \(t\) is
\(\sqrt{8t}\ll\hbar c/(\hbar c\Lambda)\), where the flowed correlators
are the free ones up to \(O(g^2(\sqrt t))\). There the bound reads
\(m\lesssim4.26\,\hbar c/\sqrt{8t}\), which is weaker than
\(\hbar c\Lambda\) by the ratio \(1/(\sqrt{8t}\Lambda)\gg1\): no
constraint on the gap. The bound reaches the expected value of \(m\)
only at \[\sqrt{8t_*}\ \simeq\ 4.26\,\frac{\hbar c}{m},\] the
correlation length, and there the coefficient is a nonperturbative
number, since the leading behaviour of \(\chi_{\varphi_t}\) at
\(\sqrt{8t}\sim\hbar c/m\) is not given by the Gaussian formula of
Section 3. This sharpens the statement of
\href{finiteness-half-flowed-susceptibility.md}{the finiteness note} §5:
the breakdown of the \(1/\sqrt{8t}\) scaling carries the content of the
problem. \textbf{Proving \(m<\infty\) requires the value of one
susceptibility at the confinement scale, and no expansion around the
free theory supplies it.}

\emph{Which susceptibility.} At \(\sqrt{8t}\sim\hbar c/m\),
\(\chi_{\varphi_t}\) is the zero-momentum fluctuation of the energy
density smeared over a correlation length, which in a gapped theory is
finite and positive, and in a theory with no finite-energy excitations
would have to vanish. It is therefore the natural target for a
nonperturbative lower bound, and a lattice statement of the form
\(\chi_{\varphi_t}\ge c\,t^{-5/2}\) uniform in the cutoff would close
the finiteness half.

\hypertarget{consequence-for-state}{%
\subsection{8. Consequence for STATE}\label{consequence-for-state}}

The finiteness half is now quantitative: an exact inequality, a closed
free-field value with an independent check, and an explicit statement of
where it becomes informative. The next step it names is a lower bound on
\(\chi_{\varphi_t}\) at \(\sqrt{8t}\) of order the correlation length,
uniform in the lattice spacing, which is a positivity statement about
one flowed correlator rather than a spectral statement. Together with
hypothesis (F1) it would give \(m<\infty\); the clause \(m>0\) remains
untouched by this whole line, and its obstruction is the one recorded in
\href{schur-error-ultraviolet.md}{the Schur note}.
\end{document}
